\section{Búsqueda y selección de múltiples candidatos}

\subsection{Muestreo autoconsistente y candidatos múltiples}
En el nivel de search, la autoconsistencia (Self-Consistency) se convierte en el arma central: permite que el modelo genere varios candidatos y después usa el voto por mayoría o el consistency score para elegir la final answer más estable. Como se aprecia en la figura \ref{fig:self-consistency}, a medida que el número de candidatos sube de 10 a 40, la accuracy sigue mejorando, muy por encima del ordenamiento basado en log-prob.

\begin{figure}[H]
\centering
\includegraphics[width=0.92\textwidth]{frame-self-consistency.jpg}
\caption{Curvas de Self-Consistency que comparan el ordenamiento log-prob con otras estrategias de sampling.\protect\footnotemark}
\label{fig:self-consistency}
\end{figure}
\footnotetext{Intervalo de tiempo en el video: 00:45:05--00:50:05; en la figura, el muestreo diversificado y la votación superan claramente al ordenamiento log-prob único.}

\begin{knowledgebox}{Tres elementos del muestreo con múltiples candidatos}
1. Diversidad: una temperature alta, el nucleus sampling y la perturbación del orden del prompt garantizan la diferencia entre los candidate; 2. Estrategia de selección: el voto por mayoría es mejor que el log-prob, salvo que se haya entrenado un verifier fuerte para cierta tarea; 3. Asignación de recursos: el multi-sample solo se activa en escenarios que necesitan una exactitud más alta y toleran la latencia.
\end{knowledgebox}

\subsection{Sinergia entre Self-Consistency y Search}
Self-Consistency le da a search un width más amplio, mientras que beam search/Tree-of-Thought controla el depth mediante stepwise scoring. Chen señala en el minuto 47:02 que el flujo ideal es primero usar sampling para generar entre 20 y 40 candidate, después usar search + verifier para decidir qué ruta puede seguir profundizándose y, al final, usar la votación de consistency para estabilizar la answer final. Usar simplemente el majority vote es confiable, pero para controlar el costo en inference-time todavía es necesario coordinarse con search: por un lado se garantiza la diversidad de los candidatos y, por otro, se usa el verifier para cortar los bad paths desde early stage.

\begin{warningbox}{No depender solo de uno de los módulos}
Si solo se ejecuta Self-Consistency sin agregar ningún search, se puede agotar el costo de token sin dejar de repetir errores similares; a la inversa, si solo se ejecuta search sin agregar sampling, la mayoría de las veces solo se alcanza un óptimo local. Por eso es necesario encadenar, dentro del inference-time pipeline, la generación de sample, la expansión de search y la votación de consistency.
\end{warningbox}

\subsection{Verificadores y líneas base de ordenamiento}
Si se entrena un stepwise verifier, se pueden filtrar las ramas de baja calidad desde el inicio de la generación; el trace de AlphaCode de DeepMind es un ejemplo de esto, pues documenta por sí misma la code execution y otorga una puntuación en cada paso, lo cual influye en el search posterior. Un verifier escalar debe equilibrar la estabilidad y la interpretabilidad, de lo contrario puede juzgar erróneamente una respuesta correcta y compleja como una «caja negra no interpretable».

\begin{warningbox}{La trampa de un verifier débil}
Si el verifier se inclina en exceso hacia las salidas cortas o hacia las plantillas del conjunto de entrenamiento, calificará el razonamiento de varios pasos como «puntuación baja» y lo descartará, lo cual hace que self-consistency converja hacia un conjunto de soluciones comunes pero incorrectas. El verifier debe ser sensible tanto a la accuracy como a la reasoning depth.
\end{warningbox}

\subsection{Pistas de ingeniería para la búsqueda}
En la diapositiva del minuto 52:50, Chen cuantifica aún más la salida de sampling, search y verifier: métricas como el conteo de muestreo, candidate score y voter margin deben revisarse continuamente en los inference-time logs. Un pipeline común es generar 30 response con high-temperature sampling, asignar una puntuación a cada response con un stepwise verifier y después usar majority vote para coronar la solución final. Cada etapa debe tener su propia tabla de monitoreo, de modo que en un despliegue real se pueda detectar el warning de «falta de diversity en el muestreo de hoy» o «verifier drift».

\begin{knowledgebox}{Métricas del pipeline de búsqueda}
1. Sampling depth: cantidad de candidatos y diversity; 2. Verifier quality: tendencia y estabilidad del stepwise score; 3. Voting stability: voter margin o consistency rate, usados para decidir si se necesita search adicional.
\end{knowledgebox}

\subsection{Resumen de la sección}
La autoconsistencia junto con un muestreo diversificado permite que el modelo se recupere de mistakes, pero debe combinarse con un ordenamiento por verifier de alta calidad. Todo el mecanismo requiere equilibrar accuracy, token budget y latency.

